% hw9-planarity-coloring.tex

% !TEX program = xelatex
%%%%%%%%%%%%%%%%%%%%
% see http://mirrors.concertpass.com/tex-archive/macros/latex/contrib/tufte-latex/sample-handout.pdf
% for how to use tufte-handout
\documentclass[a4paper, justified]{tufte-handout}

\input{hw-preamble} % feel free to modify this file if you understand LaTeX well
%%%%%%%%%%%%%%%%%%%%
\title{计算机数学 (10-planarity-coloring: 平面图与图着色)}
\me{魏恒峰}{hfwei@hnu.edu.cn}{}{}
\date{2026年5月30日}
%%%%%%%%%%%%%%%%%%%%
\begin{document}
\maketitle
%%%%%%%%%%%%%%%%%%%%
\noplagiarism % PLEASE DON'T DELETE THIS LINE!
%%%%%%%%%%%%%%%%%%%%
\begin{abstract}
\end{abstract}
%%%%%%%%%%%%%%%%%%%%
\beginrequired
%%%%%%%%%%%%%%%

%%%%%%%%%%%%%%%
\begin{problem}[平面图]
  假设 $G$ 是顶点数 $\ge 11$ 的简单图, $\overline{G}$ 是 $G$ 的补图~\footnote{
    补图: 顶点集相同, 但是 $e$ 是 $G$ 的边当且仅当 $e$ 不是 $\overline{G}$ 的边。
  }。

  \begin{enumerate}[(1)]
    \item 请证明: $G$ 和 $\overline{G}$ 不同为平面图。
    \item 请举例说明: 当 $G$ 的顶点数为 $8$ 时,
      $G$ 和 $\overline{G}$ 可以同时为平面图。
  \end{enumerate}
\end{problem}

\begin{proof}
\end{proof}
%%%%%%%%%%%%%%%

%%%%%%%%%%%%%%%
\begin{problem}[图同胚]
  假设图 $G_{1}$ 与 $G_{2}$ 是 homeomorphic (同胚) 的。
  请证明~\footnote{$m$, $n$ 分别表示边数与点数。}:
  \[
    m_{1} - n_{1} = m_{2} - n_{2}.
  \]
\end{problem}

\begin{proof}
\end{proof}
%%%%%%%%%%%%%%%

%%%%%%%%%%%%%%%
\begin{problem}[图着色]
  假设 $G$ 是包含 $n$ 个顶点的$d$-正则简单图。

  \noindent 请证明
  \[
    \chi(G) \ge \frac{n}{n-d}\text{。}
  \]
\end{problem}

\begin{proof}
\end{proof}
%%%%%%%%%%%%%%%

%%%%%%%%%%%%%%%
\begin{problem}[平面图的着色]
  假设 $G$ 是不包含三角形 $\triangle$ 的简单平面图。
  \begin{enumerate}[(1)]
    \item 请使用Euler公式证明 $G$ 含有度数 $\le 3$ 的顶点。
    \item 请使用数学归纳法证明 $G$ 是 $4$-可着色的。
  \end{enumerate}
\end{problem}

\begin{proof}
\end{proof}
%%%%%%%%%%%%%%%

%%%%%%%%%%%%%%%
\begin{problem}[着色多项式]
  设 $C_{n}$ 是 $n$ 个顶点的圈图。
  请给出 $C_{n}$ 的着色多项式 $P(C_{n}, k)$。

  \noindent 要求给出证明或计算过程。
\end{problem}

\begin{solution}
\end{solution}
%%%%%%%%%%%%%%%

%%%%%%%%%%%%%%%
\begin{problem}[图着色的应用]
  现在需要为七场讲座制定一份时间表。
  由于有些学生希望参加多场讲座，某些讲座不能安排在同一时间段。
  下表中用星号标出了不能同时进行的讲座对。

  \noindent 请问，安排这七场讲座至少需要多少个时间段？
  \fig{width = 0.50\textwidth}{figs/timetable}
\end{problem}

\begin{solution}
\end{solution}
%%%%%%%%%%%%%%%

%%%%%%%%%%%%%%%%%%%%
% 如果没有需要订正的题目，可以把这部分删掉
% \begincorrection
%%%%%%%%%%%%%%%%%%%%

%%%%%%%%%%%%%%%%%%%%
% 如果没有反馈，可以把这部分删掉
\beginfb

你可以写 (也可以发邮件)
\begin{itemize}
  \item 对课程及教师的建议与意见
  \item 教材中不理解的内容
  \item 希望深入了解的内容
  \item $\cdots$
\end{itemize}
%%%%%%%%%%%%%%%%%%%%
\end{document}